% part: intuitionistic-logic
% chapter: propositions-as-types
% section: type-preservation

\documentclass[../../../include/open-logic-section]{subfiles}

\begin{document}

\olfileid{int}{pty}{tpr}

\olsection{حفظ نوع}

تبدیل‌های کاهشی و جایگشتیِ بررسی‌شده را می‌توان مستقیماً برای
دستگاه‌های~\Log{tN2i} و~\Log{tN3i} نیز تعریف کرد. دستگاه~\Log{tN3i}
نوع ترم~$\lambd$ایِ~$N$ را نسبت به بافت~$\Gamma$، !!{formula}~$!A$
تعریف می‌کند هرگاه~$\Gamma \Sequent N : !A$ در~\Log{tN3i}
!!a{proof} داشته باشد. این واقعیت که قواعد کاهش بتای یک‌گامی برای
ترم‌های~$\lambd$ دقیقاً بازتاب تبدیل‌های کاهشی و جایگشتیِ~!!{proof}s
هستند پیامد مهمی دارد: نوع ترم~$\lambd$ نسبت به بافت~$\Gamma$ پس از
یک گام کاهش بتا ثابت می‌ماند. این ویژگی را \emph{حفظ نوع} می‌نامند.

\begin{prop}
اگر~$N$ نسبت به~$\Gamma$ نوع~$!A$ داشته باشد و~$N \redone N'$،
آنگاه~$N'$ نیز نسبت به~$\Gamma$ نوع~$!A$ دارد.
\end{prop}

\begin{proof}
  با استقرا بر~$N$ و ارتفاع~!!{proof}s~$\delta$ از
  $\Gamma \Sequent N : !A$ به‌آسانی مطلب زیر را ثابت می‌کنیم: اگر~$M$
  زیرترمی از~$N$ باشد، آنگاه~$\delta$ شامل زیر-!!{proof}ی چون
  ~$\delta_1$ است که به~$\Gamma' \Sequent M : !B$ پایان می‌یابد. اگر
  ~‌$M$ ردکس باشد، تبدیلی بر~$\delta_1$ اعمال می‌شود. با اعمال کاهش
  بر~$\delta_1$، !!a{proof}~$\delta_1'$ از
  $\Gamma' \Sequent M' : !B$ به دست می‌آوریم. با وارسی تبدیل‌های
  ~\Log{tN3i} می‌بینیم~$M'$ حاصل کاهش~$M$ است. جایگزین‌کردن
  ~$\delta_1$ در~$\delta$ با~$\delta_1'$، !!a{proof}~$\delta'$ از
  $\Gamma \Sequent N' : !A$ می‌دهد، که در آن~$N'$ از جایگزین‌کردن
  زیرترم~$M$ در~$N$ با~$M'$ حاصل شده است.

  برای نمونه، تبدیل کاهشیِ~\Intro{\land} و سپس~\Elim{\land} را در نظر
  بگیرید:
  \[
  \bottomAlignProof
  \AxiomC{}
  \Deduce$\Gamma \fCenter N_1 : !B_1$
  \AxiomC{}
  \Deduce$\Gamma \fCenter N_2 : !B_2$
  \RightLabel{\Intro{\land}}
  \BinaryInf$\Gamma \fCenter \pair{N_1}{N_2} : !B_1 \land !B_2$
  \RightLabel{\Elim{\land}[i]}
  \UnaryInf$\Gamma \fCenter \proj{i}{\pair{N_1}{N_2}} : !B_i$
  \DisplayProof
  \quad\redone\quad
  \bottomAlignProof
  \AxiomC{}
  \Deduce$\Gamma \fCenter N_i : !B_i$
  \DisplayProof
  \]
  می‌بینیم ترم برهانِ~!!{proof} سمت راست، یعنی~$\delta_1'$، همان~$N_i$
  است. این دقیقاً حاصل اعمال یک گام کاهش بتا بر ترم برهان
  ~$\proj{i}{\pair{N_1}{N_2}}$ متعلق به~!!{proof} سمت چپ، یعنی
  ~$\delta_1$، است.

  یا استنتاج~\Intro{\lif} و سپس~\Elim{\lif} و تبدیل کاهشیِ اعمال‌شده بر
  آن را در نظر بگیرید:
  \begin{multline*}
  \bottomAlignProof
  \AxiomC{$\delta_2$}
  \Deduce$x: !B_1, \Gamma \fCenter N : !B_2$
  \RightLabel{\Intro{\lif}}
  \UnaryInf$\Gamma \fCenter \lambd[\typeof{x}{!B_1}][N] : !B_1 \lif !B_2$
  \AxiomC{}
  \RightLabel{$\delta_3$}
  \Deduce$\Gamma \fCenter M : !B_1$
  \RightLabel{\Elim{\lif}}
  \BinaryInf$\Gamma \fCenter (\lambd[\typeof{x}{!B_1}][N] M) : !B_2$
  \DisplayProof
  \quad\redone\\
  \bottomAlignProof
  \AxiomC{}
  \RightLabel{$\delta_3$}
  \Deduce$\Gamma \fCenter M : !B_1$
  \RightLabel{$\Subst{\delta_2}{\delta_3}{x:!B_1}$}
  \Deduce$\Gamma \fCenter \Subst{N}{M}{x} : !B_2$
  \DisplayProof
  \end{multline*}
  باز هم ترم برهانِ~!!{proof} سمت راست دقیقاً حاصل کاهش بتای یک‌گامیِ
  ترم برهانِ~!!{proof} سمت چپ است. البته باید با دقت و با استقرا بر
  ارتفاع~$\delta_2$ بررسی کرد که جایگزین‌کردن همهٔ اصول به‌شکل
  $\Gamma' \Sequent x:!B_1$ در~$\delta_2$ با~!!{proof}~$\delta_3$ از
  $\Gamma \Sequent M:!B_1$، واقعاً !!a{proof}
  $\Subst{\delta_2}{\delta_3}{x:!B_1}$ از
  $\Gamma \Sequent \Subst{N}{M}{x} : !B_2$ به دست می‌دهد.
\end{proof}

تناظر نزدیک میان~!!{proof}s و ترم‌های~$\lambd$ پیامد دیگری نیز دارد.
خاصیت \emph{پیشرفت} را باید برای ترم‌های \emph{بسته و درست‌نوع‌دار}
بیان کرد: چنین ترمی یا یک مقدار، یعنی صورت متعارف مناسب نوعش، است یا
در یک گام کاهش بتا به ترم دیگری می‌رسد. صورت‌های متعارف بسته برای نوع
تابعی، عطفی و فصلی به‌ترتیب لامبدا، جفت و تزریق‌اند؛ برای~$\lfalse$
صورت متعارف بسته‌ای وجود ندارد. این ادعا با استقرا بر اشتقاق نوع‌دهی و
لم صورت‌های متعارف اثبات می‌شود. اگر ترم بستهٔ~$N$ مقدار نباشد، وارسی
آخرین قاعدهٔ نوع‌دهی یک ردکس~$M$ در آن آشکار می‌کند؛ زیر-!!{proof}
متناظر در برهان~\Log{tN3i} را می‌توان تبدیل کرد. اگر
$N \redone N'$ با جایگزین‌کردن~$M$ با حاصل کاهش آن به دست آید، نتیجه
!!a{proof} از~$\Gamma \Sequent N' : !A$ است.

پیشرفت همراه با حفظ نوع تضمین می‌کند کاهش بتای ترم‌های بسته و
درست‌نوع‌دار «گیر نمی‌کند»: اگر ترم مقدار نباشد، به ترم درست‌نوع‌دار
دیگری با همان نوع کاهش می‌یابد.

\end{document}
